\documentclass[11pt, oneside]{article}

% =========================
% Chinese font with ctex
% Compile with XeLaTeX
% =========================
\usepackage[UTF8, scheme=plain, fontset=none]{ctex}

% 正文字体：仿宋
\setCJKmainfont{FandolFang-Regular.otf}
% 标题字体：楷体
\newCJKfontfamily\kaishu{Kai} 

% =========================
% Page geometry
% =========================
\usepackage[
	letterpaper,
	tmargin=1in,
	bmargin=1in,
	lmargin=1.1in,
	rmargin=1.1in
]{geometry}

% =========================
% Math
% =========================
\usepackage{amsmath,amsthm,amssymb,mathrsfs,verbatim,bbm,color,graphicx,authblk}
\usepackage{booktabs,multirow,makecell,colortbl,enumitem}
\usepackage{braket}
\usepackage{caption, subcaption}
\usepackage{framed}
\usepackage{bm}

% =========================
% Tables
% =========================
\usepackage{tabularx}
\usepackage{longtable}
\usepackage{array}

% 自动换行的表格列
\newcolumntype{Y}{>{\raggedright\arraybackslash}X}

% 固定宽度、左对齐
\newcolumntype{L}[1]{%
	>{\raggedright\arraybackslash}p{#1}
}

% 固定宽度、居中
\newcolumntype{C}[1]{%
	>{\centering\arraybackslash}p{#1}
}

% =========================
% Author / affiliation
% =========================
\renewcommand\Authfont{
	\fontsize{12pt}{14pt}\selectfont
}

\renewcommand\Affilfont{
	\fontsize{12pt}{14pt}\selectfont\itshape
}

\setlength{\affilsep}{0.25em}

% =========================
% Contents / section control
% =========================
\usepackage{tocloft}
\usepackage{titlesec}

% =========================
% Hyperref
% =========================
\usepackage{hyperref}

\hypersetup{
	colorlinks=true,
	linkcolor=blue,
	filecolor=blue,
	urlcolor=blue,
	citecolor=blue,
}

% =========================
% Page number: bottom center
% =========================
\makeatletter

\renewcommand{\ps@plain}{%
	\renewcommand{\@oddhead}{}%
	\renewcommand{\@evenhead}{}%
	\renewcommand{\@oddfoot}{
		\hfil
		{\fontsize{12pt}{14pt}\selectfont\thepage}
		\hfil
	}%
	\renewcommand{\@evenfoot}{
		\hfil
		{\fontsize{12pt}{14pt}\selectfont\thepage}
		\hfil
	}%
}

\makeatother

\pagestyle{plain}

% =========================
% Main font size
% 正文 12pt，正文行距 14pt
% =========================
\AtBeginDocument{
	\fontsize{13pt}{20pt}\selectfont
}

% =========================
% Paragraph
% =========================
\setlength{\parindent}{2em}
\setlength{\parskip}{0pt}

% =========================
% Section / subsection labels
% section: I. II. III.
% subsection: 1. 2. 3.
% =========================
\renewcommand{\thesection}{\Roman{section}}
\renewcommand{\thesubsection}{\arabic{subsection}}

% =========================
% Section title
% 楷体 + 粗体 + 14pt
% 左顶格
% =========================
\titleformat{\section}
    {\normalfont\kaishu\bfseries\fontsize{16pt}{17pt}\selectfont}
    {\thesection.}
    {1.0em}
    {}

\titlespacing*{\section}
    {0pt}
    {3.2ex plus 0.5ex minus 0.2ex}
    {2.0ex plus 0.3ex}

% =========================
% Subsection title
% 楷体 + 粗体 + 12pt
% 向右缩进 2em
% =========================
\titleformat{\subsection}
    {\normalfont\kaishu\bfseries\fontsize{15pt}{17pt}\selectfont}
    {\thesubsection.}
    {1.0em}
    {}

\titlespacing*{\subsection}
    {2em}
    {2.5ex plus 0.4ex minus 0.2ex}
    {1.5ex plus 0.3ex}

% =========================
% Contents style
% =========================
\renewcommand{\contentsname}{CONTENTS}

% CONTENTS 标题
\renewcommand{\cfttoctitlefont}{
	\hspace*{1.4em}
	\fontsize{11pt}{12pt}\selectfont
	\bfseries
}

\renewcommand{\cftaftertoctitle}{}

\setlength{\cftbeforetoctitleskip}{0pt}
\setlength{\cftaftertoctitleskip}{0.65em}

% 目录无点线
\renewcommand{\cftdotsep}{\cftnodots}

% =========================
% Contents indentation
%
% Roman numeral 到 VIII / IX 较长，
% 一级目录编号宽度不能只设 2em
% =========================
\setlength{\cftsecindent}{0pt}
\setlength{\cftsecnumwidth}{3.2em}

\setlength{\cftsubsecindent}{3.2em}
\setlength{\cftsubsecnumwidth}{2.2em}

% 目录编号后加 .
\renewcommand{\cftsecaftersnum}{.}
\renewcommand{\cftsubsecaftersnum}{.}

% =========================
% Contents font
% =========================
\renewcommand{\cftsecfont}{
	\fontsize{14pt}{16pt}\selectfont
}

\renewcommand{\cftsecpagefont}{
	\fontsize{14pt}{16pt}\selectfont
}

\renewcommand{\cftsubsecfont}{
	\fontsize{14pt}{16pt}\selectfont
}

\renewcommand{\cftsubsecpagefont}{
	\fontsize{14pt}{16pt}\selectfont
}

\setlength{\cftbeforesecskip}{0.8em}
\setlength{\cftbeforesubsecskip}{0.2em}

% =========================
% Title information
% =========================
\title{
	\fontsize{15pt}{20pt}\selectfont
	\bfseries
	扭转双层石墨烯中的轨道磁化文献调研报告
}

\author[1]{Qian-Sheng Shi}

\affil[1]{
	School of Physical Science and Technology,
	ShanghaiTech University, \protect\\
	Shanghai 201210, China.
}

\newcommand{\notestartdate}{August 24, 2026}

\date{
	\fontsize{11pt}{14pt}\selectfont
	(Dated: \notestartdate -- \today)
}

% =========================
% Abstract style
% =========================
\renewcommand{\abstractname}{摘要}

\renewenvironment{abstract}
{%
	\begin{center}
		{\bfseries\fontsize{12pt}{14pt}\selectfont
		\abstractname}
	\end{center}
	\vspace{-0.3em}
	\begin{list}{}{%
		\setlength{\leftmargin}{2em}
		\setlength{\rightmargin}{2em}
	}
	\item\relax
	\fontsize{12pt}{14pt}\selectfont
}
{%
	\end{list}
}

% =========================
% Custom maketitle style
% =========================
\makeatletter

\renewcommand{\@maketitle}{%
	\begin{center}
		{\@title \par}
		\vspace{1em}
		{\@author \par}
		\vspace{0.15em}
		{\@date \par}
	\end{center}
}

\renewcommand{\maketitle}{%
	\begingroup
		\let\footnote\thanks
		\@maketitle
	\endgroup
}

\makeatother

% =========================
% Abstract
% =========================
\renewcommand{\abstractname}{摘要}


% ============================================================
% DOCUMENT
% ============================================================
\begin{document}

% ============================================================
% TITLE
% ============================================================
\maketitle

% ============================================================
% ABSTRACT
% ============================================================
\begin{abstract}

这里写整份调研报告的摘要。
建议只回答四件事情：（1）调研对象是什；（2）为什么需要进行这次调研；（3）目前最重要的文献结论是什么；（4）这些结论如何影响后续研究设计。摘要不需要逐篇列举文献，也不需要展开公式推导。

\end{abstract}

\vspace{0.5em}

% ============================================================
% CONTENTS
% ============================================================
\tableofcontents

\newpage

% ============================================================
% I. 调研范围与结论先行
% ============================================================
\section{调研范围与结论先行}

本次调研的范围是什么


重要的结论是
\begin{enumerate}
    \item {moire \kaishu 体系中抑制了自选对磁性的贡献。} 超晶格的增加对应的是什么？Jian-Peng Liu和戴希的文章当中
    
    \item {\kaishu 为什么之前没有什么工作来算\(B = 0\)时刻的轨道磁化强度}
    \item {\kaishu 磁化率（极化率）计算的困难到底在哪里？}
\end{enumerate}





% ============================================================
% II. 提出假设与问题
% ============================================================
\section{提出假设与问题}
\begin{itemize}
    \item 之前文章所提出的问题
\end{itemize}




% ------------------------------------------------------------
\begin{itemize}
    \item 工作假设
\end{itemize}



\begin{enumerate}
	\item 假设 H1：……
	\item 假设 H2：……
	\item 假设 H3：……
\end{enumerate}


% ============================================================
% III. 直接相关文献分层
% ============================================================
\section{直接相关文献分层}

这一节的目的不是逐篇详细介绍，
而是快速建立整个 literature landscape。

建议根据“与当前问题的距离”进行分层，
而不是简单按照年份排序。

% ------------------------------------------------------------
\subsection{A 层：直接回答核心问题的文献}

\begin{longtable}{
	L{2.8cm}
	L{4.3cm}
	L{3.2cm}
	L{5.0cm}
}

\caption{与当前研究问题直接相关的核心文献比较}
\label{tab:literature-summary}
\\

\toprule

\textbf{文献} & \textbf{模型和参数} & \textbf{条件} & \textbf{与当前工作的启发} \\

\midrule

\endfirsthead

\multicolumn{4}{c}{
	\tablename~\thetable\ （续）
}
\\

\toprule

\textbf{文献} & \textbf{模型和参数} & \textbf{条件} & \textbf{与当前工作的启发} \\

\midrule

\endhead

\midrule
\multicolumn{4}{r}{续下页}
\\
\endfoot

\bottomrule
\endlastfoot

% ------------------------------------------------------------
% Literature 1
% ------------------------------------------------------------
作者等，年份 & 这里写模型、Hamiltonian。
&
这里写该结论成立的具体条件：
&
而写该工作对当前项目有什么作用；
哪些结果可以直接使用；
\\
\addlinespace[0.6em]
% ------------------------------------------------------------
% Literature 2
% ------------------------------------------------------------
作者等，年份 & 内容 & 内容 & 内容 \\
\addlinespace[0.6em]
% ------------------------------------------------------------
% Literature 3
% ------------------------------------------------------------
作者等，年份 & 内容 & 内容 & 内容 \\
\addlinespace[0.6em]
% ------------------------------------------------------------
% Literature 4
% ------------------------------------------------------------
作者等，年份 & 内容 & 内容 & 内容 \\

\end{longtable}


% ------------------------------------------------------------
\subsection{B 层：理论和方法直接相关的文献}

\begin{itemize}
	\item 理论 formalism；
	\item 数值方法；
	\item 关键公式；
	\item symmetry analysis；
	\item benchmark model。
\end{itemize}


% ============================================================
% IV. 逐篇读书笔记
% ============================================================
\section{逐篇读书笔记}

这一节开始真正详细阅读核心文献。

% ------------------------------------------------------------


\noindent

{\kaishu 刘建鹏，戴希的Nature Review Physics[doi number], 论证了为什么moire体系中的自旋磁性被抑制了}，judgement.并且描述
\\
\noindent

{\kaishu 刘建鹏，戴希的Nature Review Physics[doi number], 论证了为什么moire体系中的自旋磁性被抑制了}，judgement.并且描述


% ============================================================
% V. 后续工作
% ============================================================
\section{后续工作}

\begin{enumerate}
	\item 需要继续寻找的文献方向；
	\item 需要补充的关键 reference；
	\item 目前存在矛盾、需要重新核对的文献；
	\item 需要重复阅读的核心文章。
\end{enumerate}


% ------------------------------------------------------------
\subsection{待讨论问题}

这里单独保留目前没有结论、
需要之后与导师或合作者讨论的问题。

\begin{enumerate}
	\item ……
	\item ……
	\item ……
\end{enumerate}

% ============================================================
% VI. 当前项目判断与研究定位
% ============================================================
\section{当前项目判断与研究定位}

% ============================================================
% VII. 参考文献清单
% ============================================================
\section{参考文献清单}

这一节先作为本地文献索引使用。

以后如果使用 BibTeX / Biber，
可以再将这里替换为正式 bibliography。

\end{document}